% GENERATED FILE. Edit canonical lessons, then run npm run generate:books.
% canonical-source-hash: fnv1a64:36f6a743fce01e23
% canonical-lessons: LA-C528-gens, LA-C528-propinquus, LA-C528-affinis, LA-C528-consobrinus, LA-C528-gemellus, LA-R528-review-521-528

\chapter{Relatives}
\label{ch:la-relatives}
\hlchaptermodality{528}

\begin{chapteropening}
\textbf{By the end of this chapter:} \emph{I can talk about my wider family.}

Complete the last lesson of chapter 528: I can talk about my wider family.
\end{chapteropening}

\section[gēns, gentis]{gēns, gentis — a clan, family line; a people}

\label{lesson:LA-C528-gens}



\begin{quote}
Before the new one: say the Latin for sleeplessness, then the Latin for to recline.
\end{quote}

\subsection*{You'll want to know: gēns, gentis}
\textbf{gēns, gentis} — \textquotedblleft{}a clan, family line; a people\textquotedblright{}.

\begin{quote}
Haec gēns antīqua et nōbilis est.
\end{quote}

\begin{quote}
\emph{This clan is old and noble.}
\end{quote}

Feminine, third declension. A Roman's family name showed his \emph{gēns}. English \emph{gentle} and \emph{genteel} come from it.

\begin{cousinweb}
\textbf{gēns} is inherited from a Proto-Indo-European word for \textquotedblleft{}birth, family\textquotedblright{}, built on the root of \emph{gignō} \textquotedblleft{}to beget\textquotedblright{} (EtymDB; Lewis \& Short). English \emph{gentle} and \emph{gentile} come from it, and English \emph{kin} and \emph{kind} are cousins.
\end{cousinweb}

\begin{grammarlens}[title={What you've built}]
One more word for telling, explaining and describing at B1.
\end{grammarlens}

\subsection*{Your turn}
\begin{itemize}
\raggedright
  \item \emph{Say it:} \emph{gēns, gentis}
  \item \emph{Say it:} \emph{gēns, gentis}, once more
  \item \emph{Say it:} a sentence of your own with \emph{gēns, gentis}
  \item \emph{Recall:} say the Latin for sleeplessness, then the Latin for to recline, then say \emph{gēns, gentis} again
\end{itemize}

\begin{tcolorbox}[breakable,colback=teal!4,colframe=teal!35!black,title={Before you move on}]
What does gēns, gentis mean? (\textquotedblleft{}A clan, family line; a people\textquotedblright{}.) Say it once more.
\end{tcolorbox}
\section[propinquus, propinqua, propinquum]{propinquus, propinqua, propinquum — near; a relative}

\label{lesson:LA-C528-propinquus}



\begin{quote}
Before the new one: say the Latin for to recline, then the Latin for a clan.
\end{quote}

\subsection*{You'll want to know: propinquus, propinqua, propinquum}
\textbf{propinquus, propinqua, propinquum} — \textquotedblleft{}near; a relative\textquotedblright{}.

\begin{quote}
Multī propinquī ad nūptiās veniunt.
\end{quote}

\begin{quote}
\emph{Many relatives come to the wedding.}
\end{quote}

From \emph{prope} — near. As a noun, \emph{propinquus} or \emph{propinqua} — a relative. English \emph{propinquity}.

\begin{cousinweb}
\textbf{propinquus} is built on \emph{prope} \textquotedblleft{}near\textquotedblright{} (Lewis \& Short). English \emph{propinquity} \textquotedblleft{}nearness\textquotedblright{} comes from it.
\end{cousinweb}

\begin{grammarlens}[title={What you've built}]
One more word for telling, explaining and describing at B1.
\end{grammarlens}

\subsection*{Your turn}
\begin{itemize}
\raggedright
  \item \emph{Say it:} \emph{propinquus, propinqua, propinquum}
  \item \emph{Say it:} \emph{propinquus, propinqua, propinquum}, once more
  \item \emph{Say it:} a sentence of your own with \emph{propinquus, propinqua, propinquum}
  \item \emph{Recall:} say the Latin for to recline, then the Latin for a clan, then say \emph{propinquus, propinqua, propinquum} again
\end{itemize}

\begin{tcolorbox}[breakable,colback=teal!4,colframe=teal!35!black,title={Before you move on}]
What does propinquus, propinqua, propinquum mean? (\textquotedblleft{}Near; a relative\textquotedblright{}.) Say it once more.
\end{tcolorbox}
\section[affīnis, affīne]{affīnis, affīne — related by marriage; an in-law}

\label{lesson:LA-C528-affinis}



\begin{quote}
Before the new one: say the Latin for a clan, then the Latin for near.
\end{quote}

\subsection*{You'll want to know: affīnis, affīne}
\textbf{affīnis, affīne} — \textquotedblleft{}related by marriage; an in-law\textquotedblright{}.

\begin{quote}
Socer meus affīnis meus est.
\end{quote}

\begin{quote}
\emph{My father-in-law is my relative by marriage.}
\end{quote}

Two endings. As a noun, a relative by marriage, as against a \emph{cognātus} by blood. English \emph{affinity}.

\begin{cousinweb}
\textbf{affīnis} is built on \emph{ad-} \textquotedblleft{}at\textquotedblright{} and \emph{fīnis} \textquotedblleft{}boundary\textquotedblright{} (Lewis \& Short; EtymDB): first a neighbour, then a relative by marriage. English \emph{affinity} comes from it.
\end{cousinweb}

\begin{grammarlens}[title={What you've built}]
One more word for telling, explaining and describing at B1.
\end{grammarlens}

\subsection*{Your turn}
\begin{itemize}
\raggedright
  \item \emph{Say it:} \emph{affīnis, affīne}
  \item \emph{Say it:} \emph{affīnis, affīne}, once more
  \item \emph{Say it:} a sentence of your own with \emph{affīnis, affīne}
  \item \emph{Recall:} say the Latin for a clan, then the Latin for near, then say \emph{affīnis, affīne} again
\end{itemize}

\begin{tcolorbox}[breakable,colback=teal!4,colframe=teal!35!black,title={Before you move on}]
What does affīnis, affīne mean? (\textquotedblleft{}Related by marriage; an in-law\textquotedblright{}.) Say it once more.
\end{tcolorbox}
\section[cōnsōbrīnus, cōnsōbrīnī]{cōnsōbrīnus, cōnsōbrīnī — a cousin}

\label{lesson:LA-C528-consobrinus}



\begin{quote}
Before the new one: say the Latin for near, then the Latin for related by marriage.
\end{quote}

\subsection*{You'll want to know: cōnsōbrīnus, cōnsōbrīnī}
\textbf{cōnsōbrīnus, cōnsōbrīnī} — \textquotedblleft{}a cousin\textquotedblright{}.

\begin{quote}
Cōnsōbrīnus meus in vīllā avī habitat.
\end{quote}

\begin{quote}
\emph{My cousin lives in grandfather's country house.}
\end{quote}

Masculine; \emph{cōnsōbrīna} is a female cousin. First the child of a mother's sister. English \emph{cousin} comes from it through French.

\begin{cousinweb}
\textbf{cōnsōbrīnus} is built on \emph{con-} and \emph{sōbrīnus} \textquotedblleft{}maternal cousin\textquotedblright{}, from \emph{soror} \textquotedblleft{}sister\textquotedblright{} (EtymDB; Lewis \& Short). English \emph{cousin} comes from it through Old French.
\end{cousinweb}

\begin{grammarlens}[title={What you've built}]
One more word for telling, explaining and describing at B1.
\end{grammarlens}

\subsection*{Your turn}
\begin{itemize}
\raggedright
  \item \emph{Say it:} \emph{cōnsōbrīnus, cōnsōbrīnī}
  \item \emph{Say it:} \emph{cōnsōbrīnus, cōnsōbrīnī}, once more
  \item \emph{Say it:} a sentence of your own with \emph{cōnsōbrīnus, cōnsōbrīnī}
  \item \emph{Recall:} say the Latin for near, then the Latin for related by marriage, then say \emph{cōnsōbrīnus, cōnsōbrīnī} again
\end{itemize}

\begin{tcolorbox}[breakable,colback=teal!4,colframe=teal!35!black,title={Before you move on}]
What does cōnsōbrīnus, cōnsōbrīnī mean? (\textquotedblleft{}A cousin\textquotedblright{}.) Say it once more.
\end{tcolorbox}
\section[gemellus, gemella, gemellum]{gemellus, gemella, gemellum — twin}

\label{lesson:LA-C528-gemellus}



\begin{quote}
Before the new one: say the Latin for related by marriage, then the Latin for a cousin.
\end{quote}

\subsection*{You'll want to know: gemellus, gemella, gemellum}
\textbf{gemellus, gemella, gemellum} — \textquotedblleft{}twin\textquotedblright{}.

\begin{quote}
Frātrēs gemellī simillimī sunt.
\end{quote}

\begin{quote}
\emph{The twin brothers are very alike.}
\end{quote}

A diminutive of \emph{geminus}; as a noun, \emph{gemellus} — a twin.

\begin{cousinweb}
\textbf{gemellus} is a diminutive of \emph{geminus} \textquotedblleft{}twin\textquotedblright{} (Lewis \& Short; EtymDB).
\end{cousinweb}

\begin{grammarlens}[title={What you've built}]
One more word for telling, explaining and describing at B1.
\end{grammarlens}

\subsection*{Your turn}
\begin{itemize}
\raggedright
  \item \emph{Say it:} \emph{gemellus, gemella, gemellum}
  \item \emph{Say it:} \emph{gemellus, gemella, gemellum}, once more
  \item \emph{Say it:} a sentence of your own with \emph{gemellus, gemella, gemellum}
  \item \emph{Recall:} say the Latin for related by marriage, then the Latin for a cousin, then say \emph{gemellus, gemella, gemellum} again
\end{itemize}

\begin{tcolorbox}[breakable,colback=teal!4,colframe=teal!35!black,title={Before you move on}]
What does gemellus, gemella, gemellum mean? (\textquotedblleft{}Twin\textquotedblright{}.) Say it once more.
\end{tcolorbox}
\section[(dialogue)]{Words from chapters 521–528 — a review}

\label{lesson:LA-R528-review-521-528}



\begin{quote}
Say the Latin for sick and to be ill.
\end{quote}

\subsection*{Your turn}
Say these aloud:

\begin{itemize}
\raggedright
  \item sick, to be ill, faint, weak, weakness
  \item to be pale, pale, to be red, thin, leanness
  \item sturdy, fat, slender, health, wholesome
  \item plague, an epidemic, infection, a sore, to swell
  \item to cough, to vomit, to be thirsty, bleary-eyed, inflammation of the eyes
  \item to heal, to remedy, to recover, to refresh, a draught
  \item to rest, rest, to doze, sleeplessness, to recline
  \item a clan, near, related by marriage, a cousin, twin
\end{itemize}

\begin{tcolorbox}[breakable,colback=teal!4,colframe=teal!35!black,title={Before you move on}]
Sick? (\textbf{aegrōtus, aegrōta, aegrōtum}.) To be ill? (\textbf{aegrōtō, aegrōtāre}.) Faint? (\textbf{languidus, languida, languidum}.) Weak? (\textbf{dēbilis, dēbile}.) Weakness? (\textbf{dēbilitās, dēbilitātis}.) To be pale? (\textbf{palleō, pallēre}.) Pale? (\textbf{pallidus, pallida, pallidum}.) To be red? (\textbf{rubeō, rubēre}.) Thin? (\textbf{macer, macra, macrum}.) Leanness? (\textbf{maciēs, maciēī}.) Sturdy? (\textbf{rōbustus, rōbusta, rōbustum}.) Fat? (\textbf{obēsus, obēsa, obēsum}.) Slender? (\textbf{gracilis, gracile}.) Health? (\textbf{sānitās, sānitātis}.) Wholesome? (\textbf{salūbris, salūbre}.) Plague? (\textbf{pestis, pestis}.) An epidemic? (\textbf{pestilentia, pestilentiae}.) Infection? (\textbf{contāgiō, contāgiōnis}.) A sore? (\textbf{ulcus, ulceris}.) To swell? (\textbf{tumeō, tumēre}.) To cough? (\textbf{tussiō, tussīre}.) To vomit? (\textbf{vomō, vomere}.) To be thirsty? (\textbf{sitiō, sitīre}.) Bleary-eyed? (\textbf{lippus, lippa, lippum}.) Inflammation of the eyes? (\textbf{lippitūdō, lippitūdinis}.) To heal? (\textbf{sānō, sānāre}.) To remedy? (\textbf{medeor, medērī}.) To recover? (\textbf{convalēscō, convalēscere}.) To refresh? (\textbf{recreō, recreāre}.) A draught? (\textbf{pōtiō, pōtiōnis}.) To rest? (\textbf{requiēscō, requiēscere}.) Rest? (\textbf{requiēs, requiētis}.) To doze? (\textbf{dormītō, dormītāre}.) Sleeplessness? (\textbf{īnsomnia, īnsomniae}.) To recline? (\textbf{recumbō, recumbere}.) A clan? (\textbf{gēns, gentis}.) Near? (\textbf{propinquus, propinqua, propinquum}.) Related by marriage? (\textbf{affīnis, affīne}.) A cousin? (\textbf{cōnsōbrīnus, cōnsōbrīnī}.) Twin? (\textbf{gemellus, gemella, gemellum}.)
\end{tcolorbox}
